\documentclass[crop,border=7pt]{standalone}

\usepackage{graphicx}
\usepackage{float}
\usepackage{amsmath}       
\usepackage{scalefnt}
\usepackage[dvipsnames]{xcolor}
%% TIKZ
\usepackage{tikz}
\usepackage{pgfplots}
\pgfplotsset{compat=newest}
\usetikzlibrary{decorations.pathmorphing,decorations.text,decorations.pathreplacing}
\usetikzlibrary{patterns}
\usetikzlibrary{shapes,shapes.multipart,arrows,fit,calc,positioning,intersections,through}
\usetikzlibrary{backgrounds,fadings}

%%%%%%%%%%%%%%%%%%%%%%%%%%%%%%%%%%%%%%%%%%%%%%%%%%%%%%%%%%%%%%%%%%%%%%%%%%%%%%%%%%%%%%%%%%%%%%%%%%
\begin{document}
 
 
 \tikzstyle{rect1} = [draw=red, ultra thick, rectangle, rounded corners, fill=red!20, text width=7 em, text centered, font=\bfseries]
 
 \tikzstyle{rect2} = [draw=blue, ultra thick, rectangle, rounded corners, fill=blue!20, text width=7 em, text centered, font=\bfseries]
 
 \tikzstyle{rect3} = [draw=green, ultra thick, rectangle, rounded corners, fill=green!20, text width=7 em, text centered, font=\bfseries]
 
 
 \tikzstyle{elli} = [draw=black, thick, ellipse, fill= white!20, minimum height=2em]
\tikzstyle{circ} = [draw=black, thick, circle, fill=white!20, text width=3em, text centered]
\tikzstyle{diam} = [draw=black, thick, diamond, fill=white!20, text width=4.5em, text badly centered, inner sep=0pt]
\tikzstyle{trap} = [draw=black, thick, trapezium, trapezium left angle=70, trapezium right angle=110, fill=white!20, text width=6em, text centered, minimum height=2em]
\tikzstyle{arrow} = [ultra thick,->,>=stealth]


\tikzstyle{rect4} = [draw=black, ultra thick, rectangle, rounded corners, text width=7 em, text centered, font=\bfseries]
 
\begin{tikzpicture}[>=latex, node distance=3cm]

\node (uss) [rect4] {Calcolo di autoschermaggio};
\node (track) [rect4, above of=uss, xshift=2cm] {Preparazione della matrice per il metodo delle caratteristiche};

\node (lib)  [rect4, above of=uss, xshift=-2cm] {Estrazione delle sezioni d'urto microscopiche per i diversi elementi};
\node (jeff) [rect4, above of=lib] {Librerie dei dati nucleari};
\node (geo) [rect4, above of=track] {Inizializzazione geometria assembly};
\node (asmflu) [rect4, below of=uss] {Risoluzione del sistema multigruppo per la probabilità di collisione};
\node (edi) [rect4, below of=asmflu] {Omogeneizzazione delle regioni e condesazione energetica};
\node (evo) [rect4, below of=edi] {Ottimizzazione dei calcoli di bruciamento};
\node (compo) [rect4, below of=evo] {Salvataggio delle sezioni d'urto macroscopiche ottenute};

\draw [arrow] (geo) -- (track);
\draw [arrow] (jeff) -- (lib);
\draw [arrow] (uss) -- (asmflu);
\draw [arrow] (asmflu) -- (edi);
\draw [arrow] (edi) -- (evo);
\draw [arrow] (evo) -- (compo);
\draw [arrow] (track) |- (uss);
\draw [arrow] (lib) |- (uss);
% \node (geo) [rect4, text width=8 em] {Inizializzazione della geometria 3D complessiva del reattore};
% 
% \node (macro) [rect2, below of=ini,  text width=8 em] {Interpolazione sulle sezioni d'urto e assemblaggio reattore};
% 
% \node (pow) [rect2, below of=macro] {Risoluzione del problema neutronico};
% 
% \node (flux) [rect2, below of=pow, xshift=-2cm] {Distribuzione del flusso neutronico};
% \node (power) [rect2, below of=pow,  xshift=2cm] {Distribuzione della potenza termica};
% 
% \node (pot) [rect3, right of=power, xshift=1cm] {Assegnazione sorgente termica negli assembly di combustibile};
% \node (femus) [rect3, above of=pot] {Risoluzione del problema termodinamico};
% 
% 
% \node (temp) [rect3, above of=femus] {Distribuzione di Temperatura};
% % \node (compo1) [trap, right of=dragon] {sezioni d'urto del fissile interno};
% 
% 
% \node (dragon) [rect1, left of=ini, xshift=-1cm] {Creazione database reattoristico};
% 
% \node (compo) [rect1, below of=dragon] {$\Sigma^T$ omogeneizzate e condensate};
% 
% \node (librerie) [rect1, draw=orange, fill=orange!20, above of=dragon] {Librerie di dati nucleari};
% 
% 
% \draw [arrow,orange] (librerie) -- (dragon);
% \draw [arrow, red] (dragon) -- (compo);
% \draw [arrow, red] (compo) -- (macro);
% \draw [arrow, blue] (ini) -- (macro);
% \draw [arrow, blue] (macro) -- (pow);
% \draw [arrow, blue] (pow) |- (flux);
% \draw [arrow,blue] (pow) |- (power);
% 
% \draw [arrow, green] (power) -- (pot);
% \draw [arrow, green] (pot) -- (femus);
% \draw [arrow, green] (femus) -- (temp);
% \draw [arrow, green] (temp) -- (macro);


\end{tikzpicture}
\end{document}
% pdflatex --shell-escape figure.tex
